\documentclass[preprint,1p]{elsarticle}
% \documentclass{elsarticle}
% \documentclass{article}
% \documentclass[review,leqno]{siamart1116}
%\documentclass[review,onefignum,onetabnum]{siamart190516} 



\usepackage{amsmath}
\usepackage{amssymb}
\usepackage{graphicx}
\graphicspath{{Figures/}}
\usepackage{xspace}
\usepackage{multirow}
% \usepackage{bbm}
\usepackage{mathbbol}
% mathbbol

\usepackage{algorithm}
\usepackage{algpseudocode}
% \usepackage[numbers]{natbib}
\usepackage{graphicx}
\usepackage{subcaption}
\usepackage{enumitem}

% \setlist{leftmargin=23pt}
% \setlist{labelsep=0mm}
% \setlist{rightmargin}
% \setlist{listparindent=-10pt}
% \setlist{labelwidth=10pt}
% \setlist{itemindent=0pt}

\setlist{topsep=12pt}
% \setlist{partopsep}
% \setlist{parsep=0pt}
\setlist{itemsep=-3pt}











%\usepackage{float}
% \newsiamthm{example}{Example}
\usepackage{multirow}
\usepackage{adjustbox}
\usepackage{amsmath}

\usepackage{accents}

\usepackage{Macros/macros}

% Indentation control
\setlength{\parindent}{10pt}

% Space between paragraphs.
% \setlength{\parskip}{-3pt}

%====================To be removed============================
% \usepackage[notref,notcite]{showkeys} % To see crossreferences.
% \usepackage[notref]{showkeys} % To see crossreferences.
\newcommand{\murtazo}[1]{\textcolor{blue}{#1}}
\newcommand{\mn}[1]{{\leavevmode\color{red}Murtazo: #1}}
\newcommand{\rf}[1]{{\leavevmode\color{orange}Rafa: #1}}
\usepackage[colorinlistoftodos]{todonotes}
%====================To be removed============================

\begin{document}

\begin{frontmatter}
\title{A structure-preserving scheme for an extended MHD model}
%   \tnotetext[t1]{This research is funded by Swedish Research Council (VR)
%   under grant number 2021-04620.}
  \author[1]{Murtazo Nazarov}
  \ead{murtazo.nazarov@it.uu.se}
  \author[1]{Rafael Rodriguez-Velasco}
  \ead{rafael.rodriguez-velasco@it.uu.se}
  \author[2]{Ignacio Tomas}
  \ead{igtomas@ttu.edu}
  \author[3,4]{John Shadid}
  \ead{jnshadi@sandia.gov}
  \author[3]{Michael M. Crockatt}
  \ead{mmcrock@sandia.gov}
  \address[1]{Department of Information Technology, Uppsala University, Sweden}
  \address[2]{Department of Mathematics \& Statistics Texas Tech University}
  \address[3]{Sandia National Laboratories, P.O. Box 5800, MS 1320,
              Albuquerque, NM 87185, USA}
  \address[4]{Department of Mathematics and Statistics, University of New
  Mexico, MSC01 1115, Albuquerque, NM 87131, United States}
\begin{abstract}

% We extend a previously developed structure-preserving method for the
% compressible ideal Magnetohydrodynamics (MHD) equations to the Hall resistive
% MHD model. The formulation is presented in a non-divergence formulation of the
% model, consisting of a compressible Euler’s equation solver and a
% momentum-magnetic source system with a mechanical energy update. The magnetic
% field is discretized in a curl-conforming finite element space, ensuring the
% involution constraint. The scheme preserves positivity of density, internal
% energy, conservation of total energy, and the minimum principle for specific
% entropy. We assess the accuracy of the method using whistler waves and
% present reference computations of the GEM Magnetic Reconnection Challenge and
% Orszag--Tang vortex. The method uses residual-based resistivity,
% which allows us to show stability of the scheme across various mesh configurations
% while preserving the qualitative dynamics and reconnection rates.
% The resulting method provides a robust and consistent structure-preserving
% framework for Hall MHD.
Blah.
\end{abstract}

\begin{keyword}
extended MHD \sep
structure preserving \sep
energy-stability \sep
invariant domain \sep
involution constraint
\end{keyword}

\end{frontmatter}

\section{Introduction}\label{sec:intro}

In this work we will consider the numerical solution of the
extended-MHD model:
\begin{subequations}
\label{ExtMHD}
% \left\{
\begin{align}
\label{ExtMHDmass}
\tfrac{\partial \rho}{\partial t} + \diver{}\mom &= 0 \\
%
\label{ExtMHDmom}
\tfrac{\partial \mom}{\partial t} + \diver{}
\big( \rho^{-1} \mom\mom^\transp
+ \mathbb{I} (p_i + p_e) \big)
&=
\mu \current \times \Hfield \\
%
\label{ExtMHDTotme}
\tfrac{\partial \totme}{\partial t} + \diver{}\big( \tfrac{\mom}{\rho} (\totme +
p_i + p_e) \big)
&= \Efield \cdot\current \\
%
%
\label{ExtMHDJ}
\begin{split}
\mu \tfrac{\partial \current}{\partial t}
&= \tfrac{\mu}{d_e^2} \rho( \Efield + \tfrac{\mom}{\rho} \times \Hfield) \\
%
& \ \ \ - \tfrac{\mu d_i}{d_e^2} \current \times \Hfield
- \tfrac{\mu}{d_e^2} \rho \resist \current
\end{split} \\
%
\label{ExtMHDmagneto}
\current &= \curl{}\Hfield    \, , \\
%
\label{ExtMHDHfield}
\mu \partial_t \Hfield  &= - \curl{}\Efield .
\end{align}
% \right.
\end{subequations}
where $\rho$ is the density, $\mom$ is the momentum, $\mathbb{I} \in
\mathbb{R}^{3 \times 3}$ is the identity matrix, $\totme$ is the total
mechanical energy, and $\Hfield$ is the magnetic field. Here the constant
$\resist > 0$ is the resistivity while
\begin{align}
\label{dediDef}
d_e^2 := - \tfrac{m_e m_i}{q_e q_i} > 0
  \ \ \text{and} \ \
d_i := \big(\tfrac{m_e}{q_e} + \tfrac{m_i}{q_i}\big) > 0
\end{align}
where $m_e > 0 $ is the specific electron mass, $q_e < 0 $ is the specific
electron charge, $m_i > 0 $ is the specific ion mass, $q_i > 0$ is the
specific ion charge, see for instance \cite[p. 90]{Krall1974}. The term
$\curl{}\Hfield \times \Hfield$ in the right hand side of \eqref{ExtMHDJ}
is the Hall term, while the term $\rho \resist \curl{}\Hfield$ is the
resistive term. Equation \eqref{ExtMHDmass} represents the conservation of
mass, \eqref{ExtMHDmom} is the balance of momentum, \eqref{ExtMHDTotme} is the
balance of mechanical energy, \eqref{ExtMHDJ} is commonly refereed as
generalized Ohm's law,  \eqref{ExtMHDmagneto} is the magnetostatic
condition\footnote{The magnetostatic condition is formally obtained as the
infinite speed of light of the Ampere's equation $\epsilon \partial_t \Efield -
\curl{}\Hfield = - \current$}, and \eqref{ExtMHDHfield} is the induction
equation (Faraday's law).

\begin{remark}[Ohm's law does not imply an involution]\label{Rem:NoInvolution}
Let $\boldsymbol{w}$ and $\boldsymbol{z}$ two vector fields such that
$\partial_t \boldsymbol{w} + \curl{}\boldsymbol{z} = 0$, we say that
$\boldsymbol{w}$ satisfies an divergence involution law. By this we mean that,
the evolution for $\boldsymbol{w}$ automatically implies time invariance of the
divergence. More precisely, taking the divergence to both sides of $\partial_t
\boldsymbol{w} + \curl{}\boldsymbol{z} = 0$ we obtain
\begin{align*}
\partial_t \diver{}\boldsymbol{w}
+ \diver{}\curl{}\boldsymbol{z} = \partial_t \diver{}\boldsymbol{w} = 0 \ \
\Rightarrow  \ \ \partial_t \diver{}\boldsymbol{w} = 0
\end{align*}
In other words, the divergence of $\boldsymbol{w}$ does not evolve in time.
Therefore, if the initial value of the divergence is zero, then it is zero for
all time. For instance, we say that the $\Hfield$ automatically satisfies an
involution since its evolution law is given by \eqref{ExtMHDHfield}.

On the other hand, we note that the right hand side of \eqref{ExtMHDJ}
cannot be written as a curl of some other field. In other words: Ohm's law
\eqref{ExtMHDJ} neither implies nor guarantees a divergence involution of the
current $\current$. To the best of our knowledge there: are no mathematically
nor physically meaningful arguments to claim that Ohm leads to a divergence free
condition of the current. This is a well-known issue discussed in several
references. However, it is still true that $\diver{}\current = 0$ in the
context of model \eqref{ExtMHD}. But epistemically, it is important to
understand why. In simple terms, the condition $\diver{}\current = 0$ is true
because the of the magnetostatic condition \eqref{ExtMHDmagneto}.


% More specifically:
% equation \eqref{ExtMHDJ} neither guarantees nor implies the condition
% $\diver{}\current = 0$.



% divergence
% free condition $\diver{}\boldsymbol{w}$. In more pedantic terms: the fact that
% $\partial_t
% \boldsymbol{w}$ is exactly equal to the $\curl{}$ of some other vector field
% implies a di

% arbitrary vector field $\boldsymbol{w}$ satisfying an
% evolution law of the form
% $\partial_t \boldsymbol{w} + \curl{}\boldsymbol{z} = 0$, where $\boldsymbol{z}$
% may be another given vector field
\end{remark}

% \section{Putting the current work in context}
%
% It is impossible to describe the current landscape of MHD models
% and their derivation. Indeed some aspects of their derivation are still
% controversial. But it is pedagogical to revisit the assumptions and
% the intended effect behind such assumption:
% \begin{itemize}
% \item[-] \textit{Quasi-neutrality. }The first and most important assumption of
% MHD modeling is quasi-neutrality. I wouldn't be exaggerated to say
% that \emph{quasi-neutrality assumption} and \emph{MHD-assumption} are
% colloquially used as exchangeable synonyms of each other. Quasi-neutrality is
% the foundational tenet that defines what is an MHD model and what is not.
% Quasi-neutrality does not imply the lack of electric currents: it implies the
% absence of electric charge separation. Perhaps, the most important consequence
% of quasi-neutrality is the divergence free condition $\diver{}\current = 0$. The
% original (historical) motivations to introduce the quasi-neutrality assumption
% was the removal of displacement currents and high-frequency electrostatic
% plasma-oscillations from the model. This makes MHD models (a.k.a.
% quasi-neutral models) more tractable than fully-electrodynamic two-fluid
% models.
% %
% \item[-] \textit{Infinite speed of light vs finite speed of light.} Not all
% extended MHD models use infinite speed of light in the current literature
% landscape. MhD models assuming infinite speed of light necessarily use the
% magnetostatic condition \eqref{ExtMHDmagneto} by construction. On the other
% hand, quasi-neutral models that assume finite speed of light use the whole
% Ampere's equation $\epsilon \partial_t \Efield - \curl{}\Hfield = - \current$.
% The primary motivation to re-introduce the $\epsilon > 0$ parameter is making
% $\Efield$ an evolutionary variable. This enables the re-use of many
% shock-hydrodynamics numerical schemes and code-bases. In addition, this also
% helps in avoiding global linear/nonlinear solvers. The value of $\epsilon$ used
% in practice is different from the electric permittivity of vacuum. More
% precisely we have that in practice $\epsilon > > \epsilon_0$. Or
% equivalently: in the context of quasi-neutral models $\epsilon > 0$ is a purely
% numerical device that allows to improve the computational capabilities of the
% model. Computing a transient electric field is much easier that computing  a
% quasi-static electric field that requires the solution of global linear systems.
%
% % In some cases the user might consider using the actually
% % physical value of the speed of light, but most frequently the speed of light
% % takes a purely artificial value.
% %
% \item[-] \textit{Self-consistent vs non-self-consistent models}. An MHD model
% is said to be self-consistent if $\diver{}\current = 0$ holds true by
% construction. Models that use the magnetostatic condition are self-consistent by
% default since \eqref{ExtMHDmagneto} implies:
% \begin{align*}
% \diver{}\current = \diver{}\curl{}\Hfield = 0
% \end{align*}
% On the other hand MHD models using a full Ampere's equation $\epsilon
% \partial_t \Efield - \curl{}\Hfield = 0$ are not self-consistent since Ohm's
% law does not (in general) imply and an divergence involution for
% $\current$, see Remark \ref{Rem:NoInvolution} for more details. Again, we
% highlight that this is a fairly well-known problem.
% %
% % For the very
% % specific case of the Ohm's law \eqref{ExtMHDJ} that would imply the condition:
% % \begin{align*}
% % \tfrac{1}{d_e^2} \rho( \Efield + \tfrac{\mom}{\rho} \times \Hfield)
% % - \tfrac{d_i}{d_e^2} \current \times \Hfield
% % - \tfrac{1}{d_e^2} \rho \resist \current
% % \end{align*}
% % which is a far-fetched request based on no concrete mathematical or physical
% % argument.
% %
% \item[-] \textit{Positive electron mass vs zero electron mass limit.} The most
% complete version of Ohm's law is given by:
% \begin{align*}
% bklaha.k
% \end{align*}
%
%
%
%
%
% Even for
% the lightest atomic species, which is hydrogen, we have that $m_i \approx
% 1836 m_e$. For higher atomic species (namely other gases or even metals in
% gaseous state) the ratio of ion mass to electron mass is even higher. It is
% a fair claim that the electron mass is negligible in comparison to ion mass.
% Then again, the family of quasi-neutral models available in the literature is
% divided in two: models that assume $m_e\rightarrow 0^+$, and models that
% consider $m_e > 0$.
% Ideal MHD and Resisitive-Hall-MHD assume $m_e = 0$. We often say that models
% that assume that assume zero electron mass are models that neglect all
% \emph{electron inertia effects}. On the other hand, models that assume $m_e >
% 0$ correspond with the family of quasi-neutral models often called
% \emph{extended} MHD models\footnote{We highlight that the terminology
% ``extended'' MHD model might have other uses in the literature.}
%
%
% % If we consider
% % the effect of taking the limit $m_e\rightarrow 0^+$ and then we obtain:
% % \begin{align*}
% % d_e^2 \rightarrow 0^+ \ \ \text{and} \
% % d_i \rightarrow \tfrac{m_i}{q_i}
% % \end{align*}
%
%
%
% % Applying these limit to Ohm's law \eqref{ExtMHDJ} we obtain.
% % \label{ExtMHDJ}
% % \begin{split}
% % \mu \tfrac{\partial \current}{\partial t}
% % &= \tfrac{\mu}{d_e^2} \rho( \Efield + \tfrac{\mom}{\rho} \times \Hfield) \\
% % %
% % & \ \ \ - \tfrac{\mu d_i}{d_e^2} \current \times \Hfield
% % - \tfrac{\mu}{d_e^2} \rho \resist \current
% % \end{split} \\
% \end{itemize}
%
%
% % There is a big point of departure in the MHD literature:
% % \textit{(i)}
% % models that assume infinite speed of light, therefore the magnetostatic
% % condition \eqref{ExtMHDmagneto} holds, and models that assume finite speed of
% % light, therefore
%
% % In this sense, model \eqref{ExtMHD} is self-consistent because the
% % magnetostatic condition \eqref{ExtMHDmagneto} holds. This model is similar to
% % the electrodynamic model appearing in \cite{Seyler2018}. The main difference is
% % that in \eqref{ExtMHD} we consider the infinite speed of light limit $c
% % \rightarrow \infty$ (equivalently $\epsilon \rightarrow 0$) leading to the
% % magnetostatic equation. This difference is subtle but carries a fundamental
% % difference: model \eqref{ExtMHD} is self-consistent, while model
% % \cite{Seyler2018} is not.
%
% % We note that since
% % $\tfrac{\partial
% % \current}{\partial t} \not= 0$, this model
% % includes electron kinetic effects. Therefore, the total kinetic energy is given
% % by $\tfrac{1}{2\rho}|\mom|^2 + d_e^2 \tfrac{1}{2}
% % |\frac{\current}{\sqrt{\rho}}|^2$, see for instance \RED{*** cite someone ***}.
%
%
%
%
%
%
%
%


%%%%%%%%%%%%%%%%%%%%%%%%%%%%%%%%%%%%%%%%%%%%%%%%%%%%%%%%%%%%%%%%%%%%%%%%%%%%%%%%
%%%%%%%%%%%%%%%%%%%%%%%%%%%%%%%%%%%%%%%%%%%%%%%%%%%%%%%%%%%%%%%%%%%%%%%%%%%%%%%%
%%%%%%%%%%%%%%%%%%%%%%%%%%%%%%%%%%%%%%%%%%%%%%%%%%%%%%%%%%%%%%%%%%%%%%%%%%%%%%%%





\section{Technical observation: electron
inertia regularizes the electric field}\label{EliminationOfDJDt}







This subsection is dedicated to show that model \eqref{ExtMHD} as induces a
regularization of the electric field $\Efield$. More precisely, taking the curl
of \eqref{ExtMHDHfield} we obtain:
\begin{align}
\label{CurlOfInduction}
\mu \partial_t \current = - \curl{}^2 \Efield
\end{align}
Therefore, both \eqref{ExtMHDJ} and \eqref{CurlOfInduction} provide an
evolution law for $\mu \partial_t \current$. Then if that's true, we can equate
$- \curl{}^2 \Efield$ with the right hand side of $\eqref{ExtMHDJ}$ to obtain
(after some reorganization):
\begin{align*}
\rho \Efield + \big(\tfrac{d_e^2}{\mu}\big) \curl{}^2 \Efield
= \mu d_i \current \times \Hfield
+ \tfrac{\rho}{\conduct} \current
- \mu \mom \times \Hfield
\end{align*}
which is a quasi-static definition of the electric field. Therefore we may
consider rewriting system \eqref{ExtMHD} as:
\begin{align}
\label{ExtMHDref00}
\left\{
\begin{aligned}
\tfrac{\partial \rho}{\partial t} + \diver{}\mom &= 0 \\
%
\tfrac{\partial \mom}{\partial t} + \diver{}
\big( \rho^{-1} \mom\mom^\transp
+ \mathbb{I} (p_i + p_e) \big)
&=
\mu \curl{}\Hfield \times \Hfield \\
%
\tfrac{\partial \totme}{\partial t} + \diver{}\big( \tfrac{\mom}{\rho} (\totme +
p_i + p_e) \big)
&=
\Efield \cdot\current \\
%
%
\rho \Efield + \big(\tfrac{d_e^2}{\mu}\big) \curl{}^2 \Efield
&= \mu d_i \curl{}\Hfield \times \Hfield
+ \tfrac{\rho}{\conduct} \curl{}\Hfield
- \mu \mom \times \Hfield \\
%
\mu \partial_t \Hfield  &= - \curl{}\Efield .
\end{aligned}
\right.
\end{align}
where we have decided to make explicit use of the relationship
\eqref{ExtMHDmagneto} to emphasize that the unknown fields are $\{\rho , \mom ,
\totme , \Efield,\Hfield\}$ ($\current$ is not an unknown, it is a dependent
field). We can make the description of model \eqref{ExtMHD} even more compact
if we define the differential operator:
\begin{align*}
 \mathcal{L}_{\rho} = \rho \mathbb{I}
 + \big(\tfrac{d_e^2}{\mu}\big) \curl{}^2
\end{align*}
where $\mathbb{I} \in \mathbb{R}^{3\times 3}$ is the identity operator.
Similarly, we define its solution operator as $\mathcal{L}_{\rho}^{-1}$. We
note that $\mathcal{L}_{\rho}^{-1}$ is a ``regularizing'' or ``smoothing
operator''. We may write the electric field as:
\begin{align}
\label{ExtMHDEreg}
\Efield = \mathcal{L}_{\rho}^{-1} (\mu d_i \curl{}\Hfield \times \Hfield
+ \tfrac{\rho}{\conduct} \curl{}\Hfield
- \mu \mom \times \Hfield )
\end{align}
thereby, model \eqref{ExtMHD} can be written as:
\begin{align}
\label{ExtMHDref01}
\left\{
\begin{aligned}
\tfrac{\partial \rho}{\partial t} + \diver{}\mom &= 0 \\
%
\tfrac{\partial \mom}{\partial t} + \diver{}
\big( \rho \vel\vel^\transp
+ \mathbb{I} (p_i + p_e) \big)
&=
\mu \curl{}\Hfield \times \Hfield \\
%
\tfrac{\partial \totme}{\partial t} + \diver{}\big( \vel (\totme + p_i +
p_e) \big)
&= \mathcal{L}_{\rho}^{-1} (\mu d_i \curl{}\Hfield \times \Hfield
+ \tfrac{\rho}{\conduct} \curl{}\Hfield
- \mu \mom \times \Hfield ) \cdot\current \\
%
\mu \partial_t \Hfield  &= - \curl{}\mathcal{L}_{\rho}^{-1} (\mu d_i
\curl{}\Hfield \times \Hfield
+ \tfrac{\rho}{\conduct} \curl{}\Hfield
- \mu \mom \times \Hfield ) .
\end{aligned}
\right.
\end{align}
In essence, \eqref{ExtMHDEreg} tells us that the electric field of
model \eqref{ExtMHD} is a regularization of electric
field considered in Hall-Resistive-MHD's model. Perhaps the key outtake is that
the transient nature of the current equation \eqref{ExtMHDJ} induces a
regularization of the electric field.

We note that \eqref{ExtMHDref00} and \eqref{ExtMHDref01} are just two
possible reformulations of model \eqref{ExtMHD}. However, at this point
in time, it is not evident that \eqref{ExtMHDref00} and \eqref{ExtMHDref01}
are good candidates for computational implementation. In the following section
we propose a simple reformulation of model \eqref{ExtMHD} that leads to
both a straightforward implementation and the preservation of critical
properties in the continuous, semi-discrete as well as fully-discrete setting.

%%%%%%%%%%%%%%%%%%%%%%%%%%%%%%%%%%%%%%%%%%%%%%%%%%%%%%%%%%%%%%%%%%%%%%%%%%%%%%%%
%%%%%%%%%%%%%%%%%%%%%%%%%%%%%%%%%%%%%%%%%%%%%%%%%%%%%%%%%%%%%%%%%%%%%%%%%%%%%%%%
%%%%%%%%%%%%%%%%%%%%%%%%%%%%%%%%%%%%%%%%%%%%%%%%%%%%%%%%%%%%%%%%%%%%%%%%%%%%%%%%

\subsection{A practical reformulation: eliminating the electric
field and the current}\label{EliminationOfEfield}

In this section we consider a second reformulation of model
\eqref{ExtMHD} which consists in eliminating the electric field $\Efield$
and the current $\current$. More precisely, we can eliminate the electric
field from Ohm's law \eqref{ExtMHDJ} to obtain:
\begin{align*}
\Efield = d_e^2 \tfrac{1}{\rho} \tfrac{\partial \current}{\partial t}
- \mu \vel \times \Hfield
+ \tfrac{\mu d_i}{\rho} \current \times \Hfield
+ \tfrac{1}{\conduct} \current
\end{align*}
Inserting this expression into the right hand side of
\eqref{ExtMHDTotme} and \eqref{ExtMHDHfield} and reorganizing the terms
we obtain:
\begin{align}
\label{ExtMHDNoElec}
\left\{
\begin{aligned}
\tfrac{\partial \rho}{\partial t} &= - \diver{}\mom \\
%
\tfrac{\partial \mom}{\partial t}
&=
- \diver{}
\big( \rho \vel\vel^\transp
+ \mathbb{I} (p_i + p_e) \big)
+ \mu \current \times \Hfield \\
%
\tfrac{\partial \totme}{\partial t}
- d_e^2 \tfrac{1}{\rho} \tfrac{\partial}{\partial t}\big(
\tfrac{1}{2}|\current|^2\big)
&=
- \diver{}\big( \vel (\totme + p_i +
p_e) \big)
- \mu (\vel \times \Hfield) \cdot\current
+ \tfrac{1}{\conduct} |\current|^2  \\
%
\mu \partial_t \Hfield
+ d_e^2 \curl{}( \tfrac{1}{\rho} \tfrac{\partial \current}{\partial t})
&= \curl{}(\mu \vel \times \Hfield
- \tfrac{\mu d_i}{\rho} \current \times \Hfield
- \tfrac{1}{\conduct} \current) .
\end{aligned}
\right.
\end{align}
Similarly, we substitute $\current$ with the magnetostatic condition $\current =
\curl{}\Hfield$. However, we will not make this substitution graphically
evident at all occurrences. In order to obtain compact expressions we will
stick to $\current$ when space is at a premium, and whenever the there is
enough space we will fully express its meaning by substituting with
$\curl{}\Hfield$. Therefore, we tacitly agree that there are only four
unknown fields in $\{\rho, \mom, \totme, \Hfield\}$ in formulation
\eqref{ExtMHDNoElec}. We consider the following splitting of the operator
described by \eqref{ExtMHDNoElec}:
\begin{align}
\label{SeylerOp1}
\text{Op \# 1:}
\left\{
\begin{aligned}
\tfrac{\partial \rho}{\partial t} + \diver{}\mom &= 0 \\
%
\tfrac{\partial \mom}{\partial t} + \diver{}
\big( \rho \vel\vel^\transp
+ \mathbb{I} (p_i + p_e) \big)
&= 0  \\
%
\tfrac{\partial \totme}{\partial t}
- d_e^2 \tfrac{1}{\rho} \tfrac{\partial}{\partial t}\big(
\tfrac{1}{2}|\current|^2\big)
+ \diver{}\big( \vel (\totme + p_i +
p_e) \big)
&= \bzero \\
%
\mu \partial_t \Hfield
+ d_e^2 \curl{}( \tfrac{1}{\rho} \tfrac{\partial \current}{\partial t})
&= \bzero .
\end{aligned}
\right.
\end{align}
and
\begin{align}
\label{SeylerOp2}
\text{Op \#2:}
\left\{
\begin{aligned}
\tfrac{\partial \rho}{\partial t} &= 0 \\
%
\tfrac{\partial \mom}{\partial t} &=
\mu \current \times \Hfield \\
%
\tfrac{\partial \totme}{\partial t}
- d_e^2 \tfrac{1}{\rho} \tfrac{\partial}{\partial t}\big(
\tfrac{1}{2}|\current|^2\big)
&= - \mu (\vel \times \Hfield) \cdot\current
+ \tfrac{1}{\conduct} |\current|^2  \\
%
\mu \partial_t \Hfield
+ \curl{}(d_e^2 \tfrac{1}{\rho} \tfrac{\partial \current}{\partial t})
&= \curl{}(\mu \vel \times \Hfield
- \tfrac{\mu d_i}{\rho} \current \times \Hfield
- \tfrac{1}{\conduct} \current) .
\end{aligned}
\right.
\end{align}
Here Operator \#1 has to be re-written in a more compact fashion. We start with
the following proposition that will allows us to simplify Operator \#1.

\begin{proposition}\label{PropNoDtH} Assume that that $\rho(\xcoord) >
0$ for all $\xcoord \in \domain$ and  $\Hfield \times \normal = 0$ on
$\partial\domain$. Then, the condition
\begin{align}
\label{SplitConditionH}
\mu \partial_t \Hfield
+ \curl{}(d_e^2 \tfrac{1}{\rho} \tfrac{\partial \current}{\partial t})
= \bzero \ \ \text{with} \ \ \current = \curl{}\Hfield
\end{align}
implies that $\partial_t \Hfield = 0$ and $\partial_t \current = 0$.
\end{proposition}

\begin{proof} We multiply \eqref{SplitConditionH} by $\partial_t \Hfield$ and
integrate in space to obtain:
\begin{align*}
\int_{\domain} \mu |\partial_t \Hfield|^2
+ d_e^2 \tfrac{1}{\rho}
|\partial_t\current |^2 \dx = 0 \, ,
\end{align*}
the conclusion follows readily.
\end{proof}

Therefore, from Proposition \ref{PropNoDtH}, we conclude that \eqref{SeylerOp1}
is entirely equivalent to:
\begin{align}
% \label{SeylerOp1}
\text{Op \# 1:}
\left\{
\begin{aligned}
\tfrac{\partial \rho}{\partial t} + \diver{}\mom &= 0 \\
%
\tfrac{\partial \mom}{\partial t} + \diver{}
\big( \rho \vel\vel^\transp
+ \mathbb{I} (p_i + p_e) \big)
&= 0  \\
%
\tfrac{\partial \totme}{\partial t}
+ \diver{}\big( \vel (\totme + p_i +
p_e) \big)
&= \bzero \\
%
\partial_t \Hfield &= \bzero .
\end{aligned}
\right.
\end{align}

\begin{proposition}[Properties satisfied by Operator \#2] Assume that $\Hfield
\times \normal = 0$ on the boundary $\partial\domain$. Then, the
evolution of Operator \#2, as described in \eqref{SeylerOp2} satisfies the
following properties:
\begin{align}
\label{SeylerOp2prop01}
\tfrac{\partial}{\partial t}\int_{\domain} \tfrac{1}{2}\rho |\vel|^2
+ \tfrac{\mu}{2} |\Hfield|^2
+ d_e^2 \tfrac{1}{2\rho}|\curl{}\Hfield|^2 \dx
%
&= - \int_{\domain} \tfrac{1}{\conduct} |\curl{}\Hfield|^2  \dx \\
%
\label{SeylerOp2prop02}
\tfrac{\partial }{\partial t} \big( \totme - \tfrac{1}{2}\rho
|\vel|^2 - d_e^2 \tfrac{1}{\rho} \tfrac{1}{2}|\curl{}\Hfield|^2\big)
&= \tfrac{1}{\conduct} |\curl{}\Hfield|^2 \ \ \text{for all }\xcoord \in \domain
\end{align}
Note that \eqref{SeylerOp2prop01} is a global property. However,
\eqref{SeylerOp2prop02} is a purely local property: since it does not involve
spatial gradients (e.g. divergence terms, convective terms, etc) it holds true
in a pointwise sense.
\end{proposition}

\begin{proof} We start multiplying the equation for $\partial_t\Hfield$ in
\eqref{SeylerOp2} by an arbitrary smooth test function $\Htest$ and integrate
by parts to obtain:
\begin{align*}
&\int_{\domain}\mu \partial_t \Hfield \cdot \Htest
+ d_e^2 \tfrac{1}{\rho} \tfrac{\partial \current}{\partial t}
\cdot \curl{}\Htest \dx
+ \int_{\partial\domain}
d_e^2 \tfrac{1}{\rho} \tfrac{\partial \current}{\partial t} \cdot
(\Htest \times \normal) \ds \\
%
& \ \ \ = \int_{\domain} (\mu \vel \times  \Hfield
- \tfrac{\mu d_i}{\rho} \current \times \Hfield
- \tfrac{1}{\conduct} \current) \cdot \curl{}\Htest \dx \\
%
& \ \ \ + \int_{\partial\domain} (\mu \vel \times \Hfield
- \tfrac{\mu d_i}{\rho} \current \times  \Hfield
- \tfrac{1}{\conduct} \current) \cdot
(\Htest \times \normal) \ds
\end{align*}
If we set $\Htest = \Hfield$ in the previous expression we obtain:
\begin{align}
\label{SeylerOp02formula01}
\int_{\domain} \partial_t \big(\tfrac{\mu}{2} |\Hfield|^2\big)
+ d_e^2  \partial_t \big(\tfrac{1}{2\rho} |\curl{}\Hfield|^2 \big) \dx
%
= \int_{\domain} \mu (\vel \times  \Hfield)\cdot \current
- \tfrac{1}{\conduct} |\current|^2  \dx .
\end{align}
On the other hand, multiplying the equation for $\partial_t\vel$ in
\eqref{SeylerOp2} by $\vel$ and integrating in space we obtain:
\begin{align*}
\int_{\domain} \partial_t \big(\tfrac{1}{2}\rho |\vel|^2\big) \dx &=
\int_{\domain} \mu (\current \times  \Hfield)\cdot\vel \dx \, ,
\end{align*}
adding this expression to \eqref{SeylerOp02formula01} then
\eqref{SeylerOp2prop01} follows.
\end{proof}

\begin{corollary}[Conservation of total energy] Integrating
\eqref{SeylerOp2prop02} in space:
\begin{align*}
\int_{\domain} \tfrac{\partial }{\partial t} \big( \totme - \tfrac{1}{2}\rho
|\vel|^2 - d_e^2 \tfrac{1}{2\rho} |\curl{}\Hfield|^2\big) \dx
&= \int_{\domain} \tfrac{1}{\conduct} |\curl{}\Hfield|^2 \dx
\end{align*}
adding this expression to \eqref{SeylerOp2prop01} we obtain:
\begin{align*}
\tfrac{\partial}{\partial t}\int_{\domain}
\totme + \tfrac{\mu}{2} |\Hfield|^2 \dx
%
&= 0
\end{align*}
showing that the volution of Operator \#2, as described in \eqref{SeylerOp2},
preserves the total energy.
\end{corollary}

From these results we propose the following scheme to solve Operator \#2:
\begin{itemize}
 \item[-] \textit{Step \#1.} Given the fields $\{\rho^n, \totme^{n},
\vel^{n}, \Hfield^{n}\}$ solve the coupled system:
\begin{align}
\label{SeylerOp02Scheme}
\left\{
\begin{aligned}
& \rho^n (\vel^{n+1} - \vel^{n}, \veltest) =
\dt \mu (\current^{n+ \frac{1}{2}} \times \Hfield^{n+ \frac{1}{2}}, \veltest) \\
%
& \dt \mu (\Hfield^{n+1} - \Hfield^{n}, \Htest)
+ (d_e^2 \tfrac{1}{\rho} (\curl{}\Hfield^{n+1} - \curl{}\Hfield^{n}),
\curl{}\Htest) \\
%
& \ \ \ = \dt (\mu \vel^{n+ \frac{1}{2}} \times \Hfield^{n+ \frac{1}{2}}
- \tfrac{\mu d_i}{\rho} \current^{n+ \frac{1}{2}} \times \Hfield^{n+
\frac{1}{2}}
- \tfrac{1}{\conduct} \curl{}\Hfield^{n+ \frac{1}{2}} , \curl{}\Htest )
\end{aligned}
\right.
\end{align}
\item[-] \textit{Step \#2.} With $\vel^{n+1}$ and $\Hfield^{n+1}$, now at our
disposal, we update the total mechanical energy as:
\begin{align*}
\totme^{n+1}
&=
\totme^{n}
+ \tfrac{1}{2}\rho |\vel^{n+1}|^2
- \tfrac{1}{2}\rho |\vel^{n}|^2
+ d_e^2 \tfrac{1}{2}\big|\tfrac{\curl{}\Hfield^{n+1}}{\sqrt{\rho}}\big|^2
- d_e^2 \tfrac{1}{2}\big|\tfrac{\curl{}\Hfield^{n}}{\sqrt{\rho}}\big|^2
+ \tfrac{\dt}{\conduct} |\current^{n+\frac{1}{2}}|^2
\end{align*}
 \item[-] \textit{Step \#3.} Return the updated solution
$\{\rho^n, \totme^{n+1}, \vel^{n+1}, \Hfield^{n+1}\}$. Note that the density
does not evolve during the evolution of Operator \#2 as described in
\eqref{SeylerOp2prop02}.
\end{itemize}

\begin{remark}[Well-posedness of Operator \#1] We note that the discrete
induction equation in \eqref{SeylerOp02Scheme} can reorganized as:
\begin{align*}
& \mu (\Hfield^{n+1}, \Htest)
+ ((\tfrac{d_e^2}{\rho} + \tfrac{\dt}{2 \conduct}) \curl{}\Hfield^{n+1},
\curl{}\Htest) \\
%
& + \dt \tfrac{\mu d_i}{4} (\tfrac{1}{\rho}
(\current^{n+1} \times \Hfield^{n+1} + \current^{n+1} \times \Hfield^{n}
+ \current^{n} \times \Hfield^{n+1}, \curl{}\Htest ) \\
%
& - \tfrac{\dt \mu}{4}(
\vel^{n+1} \times \Hfield^{n+1} + \vel^{n+1} \times \Hfield^{n}
+ \vel^{n} \times \Hfield^{n+1}, \curl{}\Htest) \\
%
&= \mu (\Hfield^{n}, \Htest)
+ (d_e^2 \tfrac{1}{\rho} \curl{}\Hfield^{n}, \curl{}\Htest)
+ \tfrac{\dt \mu}{4}(\vel^{n} \times \Hfield^{n}, \curl{}\Htest) \\
%
&- \dt \tfrac{\mu d_i}{4} (\tfrac{1}{\rho} \current^{n} \times \Hfield^{n},
\curl{}\Htest )
%
- \dt (\tfrac{1}{2 \conduct} \curl{}\Hfield^{n} , \curl{}\Htest ) \, ,
\end{align*}
Here the second row is related to the Hall term, while the third row is related
to the usual $\vel \times \Hfield$ term occurring in ideal in MHD. The second
row is very much benign since its spatial discretization scales like
$\tfrac{1}{h}$. However, the second row containing the Hall terms
scales like $\tfrac{1}{h^2}$. Those terms are lead to an extremely
ill-conditioned Jacobian of the Newton iteration scheme, in particular when
$\rho < < 1$. We note however that the symmetric diffusive term appearing at
the end of the first row becomes very large whenever $\rho$ is very
small, compensating the singular behaviour of the Hall term and creating (to
some extent) a better conditioned system.
\end{remark}


\RED{*** Can we add the pressure terms to the current equation? ***} \\
\RED{*** Can we add some form of convective terms that we neglected? ***} \\
\RED{*** Can we add the JJ term in the momentum equation that we neglected?
***} \\


\begin{align*}
 d_h^2 = \max \{d_e^2 , \mathcal{C} h^2\}
\end{align*}


\begin{align*}
[d_e^2] = \frac{\text{mass}^2}{\text{charge}^2} = \mathcal{C}
\cdot \text{length}^2
\end{align*}
Therefore:
\begin{align*}
\mathcal{C} = \frac{\text{mass}^2}{\text{length}^2 \text{charge}^2}
\end{align*}
We have to do either
\begin{itemize}
\item[--] Non-dimensional analysis to figure out the magnitude of  $\mathcal{C}$
\item[--] Alternatively, the value of the constant comes from a coercivity
analysis of the Jacobian.
\end{itemize}



\newpage
\section{Numerical Scheme}




\newpage
\subsection{Algorithmic Summary}




\newpage
\section{Numerical Scheme}





\newpage
\section{Computational Results}


\subsection{Solver performance statistics}


\newpage
\begin{appendix}

\section{Thermodynamics and equations of state}

\section{Non-dimensional analysis of the extended MHD system}

Consider the induction equation
\begin{align*}
    \mu \partial_t \Hfield
+ \curl{}\!\left(d_e^2 \tfrac{1}{\rho} \tfrac{\partial}{\partial t}\curl{}\Hfield\right)
&= \curl{}\!\left(\mu \vel \times \Hfield
- \tfrac{\mu d_i}{\rho} \curl{}\Hfield \times \Hfield
- \tfrac{1}{\conduct} \curl{}\Hfield\right),
\end{align*}
and the following change of variables
\begin{align*}
    \Tilde{\xcoord} = \frac{\xcoord}{L_0}, \qquad \Tilde{\vel}=\frac{\vel}{v_A}, \qquad \Tilde{\Hfield}=\frac{\Hfield}{\Hfieldcomponent_0}, \qquad \Tilde{\rho}=\frac{\rho}{\rho_0},
\end{align*}
where $v_A$ is the Alfv\'en velocity $v_A:=\Hfieldcomponent_0\sqrt{\mu/\rho_0}$ and $L_0$, $\Hfieldcomponent_0$ and $\rho_0$ are characteristic length, magnetic field strength and density, respectively. In terms of the rescaled variables, the differential operators satisfy
\begin{align*}
    \curl{} = \frac{1}{L_0}\widetilde{\curl{}}, \qquad \frac{\partial}{\partial t}=\frac{v_A}{L_0}\frac{\partial}{\partial \Tilde{t}}.
\end{align*}
Substituting these into the induction equation and dividing through by $\mu \Hfieldcomponent_0 v_A / L_0$ yields the non-dimensionalized equation
\begin{align*}
    \Tilde{\partial_t}\,\Tilde{\Hfield}
+ \alpha\,\widetilde{\curl{}}\!\left(
    \frac{1}{\Tilde{\rho}}
    \Tilde{\partial_t}\,
    \widetilde{\curl{}}\Tilde{\Hfield}
\right)
&=
\widetilde{\curl{}}\!\left(
    \Tilde{\vel}\times\Tilde{\Hfield}
\right)
-
\beta\,\widetilde{\curl{}}\!\left(
    \frac{1}{\Tilde{\rho}}
    \widetilde{\curl{}}\Tilde{\Hfield}
    \times
    \Tilde{\Hfield}
\right)
-
S^{-1}\,\widetilde{\curl{}}\!\left(
    \widetilde{\curl{}}\Tilde{\Hfield}
\right),
\end{align*}
where $S:=\mu\conduct L_0 v_A$ is the Lundquist number and the non-dimensional constants are
\begin{align*}
    \alpha = \frac{d_e^2}{L_0^2\rho_0\mu}, \qquad \beta = \frac{d_i\Hfieldcomponent_0}{L_0\rho_0 v_A}.
\end{align*}
Using the definitions of $d_i$ and $v_A$, we can write $\beta$ as
\begin{align*}
    \beta = \frac{1}{L_0}\left(\frac{m_e}{|q_e|}+\frac{m_i}{q_i}\right)\frac{1}{\sqrt{\rho_0\mu}}
    \approx\frac{1}{L_0}\frac{m_i}{q_i}\frac{1}{\sqrt{\rho_0\mu}}
    =\frac{\delta_i}{L_0}.
\end{align*}
Here we have used that $d_i\approx m_i/q_i$, and identified the ion skin depth $\delta_i:=\frac{m_i}{q_i\sqrt{\rho_0\mu}}$. The strength of the Hall term is therefore set by the ratio of the characteristic length to the ion skin depth. For numerical simulations of Hall MHD, it is common to choose $L_0=\delta_i$, so that $\beta = 1$.

We can also write $\alpha$ in terms of the ion skin depth. Substituting the definition of $d_e^2$ and multiplying and dividing by $m_i |q_e| / (q_i m_e)$:
\begin{align*}
    \alpha = \frac{1}{L_0^2}\frac{m_em_i}{\rho_0\mu |q_e| q_i}
    = \frac{1}{L_0^2}\underbrace{\frac{m_i^2}{\rho_0\mu q_i^2}}_{\delta_i^2}\,\frac{q_i m_e}{|q_e| m_i}
    = \frac{\delta_i^2}{L_0^2}\,\frac{q_i m_e}{|q_e| m_i}
    = \beta^2\,\frac{q_i m_e}{|q_e| m_i}.
\end{align*}
Hence the electron-inertia term is suppressed relative to the square of the Hall term strength by the electron-to-ion mass and charge ratio. For a hydrogen plasma, $q_i = |q_e|$ and $m_i/m_e \approx 1836$. If the characteristic length is chosen to be $L_0=\delta_i$, then $\beta = 1$ and $\alpha \approx 1/1836$.



\section{Newton solver and Jacobian analysis}





\end{appendix}


\newpage
\bibliographystyle{plain}
\bibliography{Macros/biblio}

\end{document}



